\section{Related Work}
\label{sec:related}

\paragraph{SFT for LLM agents.}
Behavioural cloning on expert trajectories bootstraps LLM agents effectively
\citep{yao2022webshop,deng2023mind2web,zeng2023agenttuning,hong2024cogagent},
but generalises poorly out-of-distribution because the agent never reasons about
downstream consequences of its actions \citep{chu2025sft}.

\paragraph{Data-collection paradigms.}
ALEE \citep{zhang2025alee} trains on next-state prediction using self-generated
rollouts, requiring live environment access plus a separate SFT stage.
ETO \citep{song2024eto} collects failure trajectories for DPO-style contrastive
learning; IPR \citep{zhang2024ipr} adds step-level reward estimation.
All three require rollouts; ETO and IPR additionally require reward signals.
\aomt{} operates on the same offline corpus as standard SFT, innovating on the
\emph{objective axis} rather than the data axis---the two are orthogonal and
composable.

\paragraph{Masked diffusion LMs.}
LLaDA \citep{nie2025llada} learns bidirectional conditional distributions
via a variational bound on token-level NLL.
LLaDA~2.0 \citep{bie2025llada20} scales this to 16B/100B parameters (MoE),
with a SFT protocol that maps directly onto \aomt{}'s context/target distinction.
XLNet \citep{yang2019xlnet} pursues any-order learning via token-level permutation,
incurring two-stream attention overhead and a train--inference discrepancy;
\aomt{} operates at the coarser, semantically meaningful trajectory-unit level.

\paragraph{Masked objectives for agents.}
STeP \citep{chen2023step} and AGENTBANK/SAMOYED \citep{zeng2023agenttuning}
use token-level masking as a loss filter on a fixed causal objective.
Neither varies the context/target assignment nor employs bidirectional context.
